% ==================================================
% PA2 REVISION HIGHLIGHTING
% ==================================================
% Requires xcolor and instructionblue in main.tex.
\providecommand{\PAChange}[1]{%
    \textcolor{instructionblue}{#1}%
}

\section{Stakeholder and User Descriptions}

This section identifies the people and organizations whose needs,
constraints, or data affect the product. Direct users are separated
from stakeholders who guide, review, or supply information to the
prototype.

\PAChange{The updated system defines five user roles:
Administrator, Patient / General User, Healthcare Professional (HCP),
Dataset Reviewer, and Technical Support. Each role has distinct
responsibilities and access permissions.}


% ==================================================
% 3.1 STAKEHOLDER SUMMARY
% ==================================================

\subsection{Stakeholder Summary}

\begin{longtable}{|
    L{0.20\textwidth}|
    L{0.27\textwidth}|
    L{0.43\textwidth}|
}
\hline
\textbf{Name / Type} &
\textbf{Description} &
\textbf{Responsibilities and Interests}\\
\hline
\endfirsthead

\hline
\textbf{Name / Type} &
\textbf{Description} &
\textbf{Responsibilities and Interests}\\
\hline
\endhead

Project Team
& Three student developers responsible for the
Medication Safety Assistant prototype.
& Define requirements, develop and test the system,
manage project progress, and ensure that the
software meets its objectives and safety requirements.
\\ \hline

Course Instructors and Teaching Assistants
& Academic supervisors and evaluators of the project.
& Review requirements, development plans,
deliverables, and project progress; provide
feedback and evaluate the final results.
\\ \hline

Potential Healthcare Advisors
& Healthcare professionals who may provide
feedback on the prototype.
& Help review the clarity and usefulness of
medication information, safety warnings, and
HCP workflows. They do not provide formal
clinical approval for the prototype.
\\ \hline

\end{longtable}


% ==================================================
% 3.2 USER SUMMARY
% ==================================================

\subsection{User Summary}

\PAChange{The application defines five user roles with
separate responsibilities. Role-Based Access Control (RBAC)
restricts each role to its permitted tasks and information.}

\begin{longtable}{|
    L{0.18\textwidth}|
    L{0.23\textwidth}|
    L{0.31\textwidth}|
    L{0.18\textwidth}|
}
\hline
\textbf{User Type} &
\textbf{Description} &
\textbf{Responsibilities / Activities} &
\textbf{Represented By}\\
\hline
\endfirsthead

\hline
\textbf{User Type} &
\textbf{Description} &
\textbf{Responsibilities / Activities} &
\textbf{Represented By}\\
\hline
\endhead

% --------------------------------------------------
% ADMINISTRATOR
% --------------------------------------------------

\PAChange{Administrator}
& \PAChange{Authorized personnel responsible
for system administration and access management.}
& \PAChange{Manage user accounts, roles, and
permissions; monitor authorized administrative
activities and system records.}
& \PAChange{Project team acting as administrators
during prototype testing.}
\\ \hline

% --------------------------------------------------
% PATIENT / GENERAL USER
% --------------------------------------------------

Patient / General User
& People who take one or more medications,
especially those managing multiple prescriptions.
& Enter medication information through OCR,
NLP, Voice, or Manual Input; review and confirm
medication details; check safety warnings;
request professional review when needed;
manage confirmed schedules and track adherence.
\PAChange{Use the AI chatbot, review medication history,
prepare health summaries, and communicate with
HCPs within authorized review cases.}
& Potential patient representatives and
planned user feedback.
\\ \hline

% --------------------------------------------------
% HEALTHCARE PROFESSIONAL
% --------------------------------------------------

\PAChange{Healthcare Professional (HCP)}
& \PAChange{A simulated professional role
representing qualified healthcare professionals
within the research prototype.}
& Review medication summaries, detected risks,
and uncertain cases; request corrections
or missing information; record review decisions
and confirm proposed schedule changes.
\PAChange{Exchange messages with patients in authorized
review cases and access permitted health summaries
and medication history.}
& Project team and potential
healthcare professional feedback.
\\ \hline

% --------------------------------------------------
% DATASET REVIEWER
% --------------------------------------------------

\PAChange{Dataset Reviewer}
& \PAChange{Authorized personnel responsible
for reviewing medication dataset quality
and supporting evidence.}
& \PAChange{Inspect medication records,
ingredient information, source references,
and mappings; identify inconsistencies
and record review decisions.}
& \PAChange{Project team acting as dataset
reviewers during prototype testing.}
\\ \hline

% --------------------------------------------------
% TECHNICAL SUPPORT
% --------------------------------------------------

\PAChange{Technical Support}
& \PAChange{Authorized personnel responsible
for assisting users with technical issues.}
& \PAChange{Investigate reported errors,
review permitted diagnostic information,
document solutions, and track issue resolution
without unnecessary access to patient data.}
& \PAChange{Project team acting as technical
support during prototype testing.}
\\ \hline

\end{longtable}


% --------------------------------------------------
% USER CHARACTERISTICS (design assumptions)
% --------------------------------------------------

\PAChange{\textbf{User Characteristics.}
The following profiles describe expected users for
interface design, including Vietnamese-speaking adults
and older patients with varied digital experience.
These are design assumptions, not results from
completed user interviews or clinical testing.}

\begin{longtable}{|
    L{0.18\textwidth}|
    L{0.22\textwidth}|
    L{0.28\textwidth}|
    L{0.22\textwidth}|
}
\hline
\textbf{User Type} &
\textbf{Expected Skills} &
\textbf{Needs and Expectations} &
\textbf{Main Constraints}
\\ \hline
\endfirsthead
\hline
\textbf{User Type} &
\textbf{Expected Skills} &
\textbf{Needs and Expectations} &
\textbf{Main Constraints}
\\ \hline
\endhead

\PAChange{Patient / General User}
& \PAChange{Basic or varied experience with smartphones and Web forms; no medical training assumed.}
& \PAChange{Clear medicine names, understandable warnings, simple confirmation, and easy reminder use.}
& \PAChange{Small screens, poor prescription images, limited digital or health literacy, and accessibility needs.}
\\ \hline

\PAChange{Healthcare Professional (HCP)}
& \PAChange{Medical knowledge is represented in a simulated account; basic Web skills.}
& \PAChange{Well-organized cases, data limits, replies, and traceable decisions.}
& \PAChange{Limited review time and access only to authorized patient information.}
\\ \hline

\PAChange{Dataset Reviewer}
& \PAChange{Able to inspect structured medication data and supporting references.}
& \PAChange{Clear source records, mapping status, and recorded review decisions.}
& \PAChange{Incomplete references, ambiguous names, and uncertain mappings.}
\\ \hline

\PAChange{Administrator}
& \PAChange{Basic account-management and access-control knowledge.}
& \PAChange{Manage accounts, roles, and view permitted activity records.}
& \PAChange{Must avoid incorrect permissions and unauthorized access.}
\\ \hline

\PAChange{Technical Support}
& \PAChange{Basic troubleshooting and Web-system knowledge.}
& \PAChange{Clear issue reports, status history, and permitted diagnostic details.}
& \PAChange{Limited access to private patient information.}
\\ \hline
\end{longtable}


% ==================================================
% 3.3 USER ENVIRONMENT
% ==================================================

\subsection{User Environment}

The product is used through modern desktop and mobile Web browsers.

\PAChange{Patients mainly use the application at home,
after receiving prescriptions or managing existing
medication schedules. HCPs, Dataset Reviewers,
Administrators, and Technical Support use restricted
role-specific Web interfaces, mainly on desktop or tablet.}

\PAChange{Camera, microphone, and notifications depend
on device and browser support and user permissions.
Manual entry remains available when OCR or voice
input is unavailable. Limited connectivity, poor
lighting, background noise, and small screens can
make input and review difficult.}

\PAChange{The interface should provide readable text,
editable extracted information, clear safety limits,
and warnings that do not depend on color alone.
The AI chatbot supports multi-turn requests, while
a separate conversation within each authorized HCP
review case supports patient--HCP replies.}

\PAChange{\textbf{User Task Analysis.} The table maps
current tasks to the 18 use cases. Frequencies are
expected design assumptions, not observed usage data.}

\begin{longtable}{|
    L{0.17\textwidth}|
    L{0.31\textwidth}|
    L{0.21\textwidth}|
    L{0.21\textwidth}|
}
\hline
\textbf{Actor} &
\textbf{Task and Goal} &
\textbf{When / Frequency} &
\textbf{Environment / Constraint}
\\ \hline
\endfirsthead
\hline
\textbf{Actor} &
\textbf{Task and Goal} &
\textbf{When / Frequency} &
\textbf{Environment / Constraint}
\\ \hline
\endhead

\PAChange{Patient}
& \PAChange{Enter and confirm medication information (UC-02); obtain a reliable record.}
& \PAChange{When receiving a new prescription or correcting captured input.}
& \PAChange{Home; camera quality, typing difficulty, or network limits.}
\\ \hline

\PAChange{Patient}
& \PAChange{Check medication safety (UC-03); understand findings and limits.}
& \PAChange{After confirming a medication list or adding a new prescription.}
& \PAChange{Mobile or desktop; incomplete medicine data.}
\\ \hline

\PAChange{Patient / HCP}
& \PAChange{Request and review a safety case, with follow-up replies (UC-04, UC-05).}
& \PAChange{When review is needed; responses may occur later.}
& \PAChange{Separate authorized sessions; simulated HCP workflow.}
\\ \hline

\PAChange{Patient}
& \PAChange{Set or manage reminders and record dose status (UC-06--UC-09).}
& \PAChange{When a verified schedule exists; dose tracking may occur daily.}
& \PAChange{Mostly mobile; notification permissions and verified schedules.}
\\ \hline

\PAChange{Patient}
& \PAChange{View records, ask medication questions, and prepare or share a health summary (UC-10--UC-13).}
& \PAChange{As needed, often before a healthcare discussion.}
& \PAChange{Home or mobile; available records and sharing permission.}
\\ \hline

\PAChange{Dataset Reviewer}
& \PAChange{Review medication dataset records and evidence (UC-14).}
& \PAChange{During dataset updates or quality review.}
& \PAChange{Desktop; incomplete or conflicting source evidence.}
\\ \hline

\PAChange{Administrator}
& \PAChange{Manage accounts, roles, and view audit history (UC-15, UC-16).}
& \PAChange{When accounts or permissions change; periodic checks.}
& \PAChange{Restricted desktop interface; RBAC rules.}
\\ \hline

\PAChange{All registered roles}
& \PAChange{Sign in (UC-01) and report technical issues (UC-17).}
& \PAChange{When accessing the system or encountering a technical problem.}
& \PAChange{Web browser; account access and issue details.}
\\ \hline

\PAChange{Technical Support}
& \PAChange{Handle reported technical issues and update support tickets (UC-18).}
& \PAChange{When a support ticket needs attention.}
& \PAChange{Restricted Web interface; access only to permitted diagnostic information.}
\\ \hline
\end{longtable}

The prototype uses synthetic or de-identified data
and is not integrated with hospital, pharmacy,
electronic health-record, or real professional-review
systems.


% ==================================================
% 3.4 SUMMARY OF KEY STAKEHOLDER OR USER NEEDS
% ==================================================

\subsection{Summary of Key Stakeholder or User Needs}

\begin{longtable}{|
    L{0.21\textwidth}|
    L{0.09\textwidth}|
    L{0.20\textwidth}|
    L{0.20\textwidth}|
    L{0.20\textwidth}|
}
\hline
\textbf{Need} &
\textbf{Priority} &
\textbf{Concerns / Cause} &
\textbf{Current Solution} &
\textbf{Proposed Solution}\\
\hline
\endfirsthead

\hline
\textbf{Need} &
\textbf{Priority} &
\textbf{Concerns / Cause} &
\textbf{Current Solution} &
\textbf{Proposed Solution}\\
\hline
\endhead

% --------------------------------------------------
% MEDICATION INTAKE
% --------------------------------------------------

Enter medication information easily
& High
& Long prescriptions, unfamiliar names,
and difficulty typing.
& Manually enter information or ask others for help.
& Support OCR, natural-language text,
voice, and manual input.
\\ \hline

Ensure accurate medication information
& High
& OCR, NLP, or voice input may produce
incorrect medication details.
& Manually check prescriptions and labels.
& Allow users to review, correct,
and confirm medication information
before safety assessment.
\\ \hline

% --------------------------------------------------
% MEDICATION IDENTIFICATION AND SAFETY
% --------------------------------------------------

Identify Vietnamese medications
& High
& Different brand names may contain
the same or related active ingredients.
& Search medicine labels or websites.
& Identify medicines and active ingredients
using supported Vietnamese medication data.
\\ \hline

Understand medication safety risks
& High
& Users may not recognize drug interactions
or duplicate active ingredients.
& Search online or consult a healthcare professional.
& Check supported drug interactions and
duplicate ingredients; provide clear
risk levels and warnings.
\\ \hline

Understand incomplete or uncertain results
& High
& Missing medication details or limited data
may lead to misleading conclusions.
& Compare information from different sources.
& Report missing information, unresolved
medicines, and assessment limitations.
Never treat unknown results as safe.
\\ \hline

% --------------------------------------------------
% HEALTHCARE PROFESSIONAL REVIEW
% --------------------------------------------------

Know when professional review is needed
& High
& Users may misunderstand safety warnings
or change medications without advice.
& Contact a doctor or pharmacist when possible.
& Provide clear warnings and identify
high-risk or uncertain cases that
require professional review.
\\ \hline

Review medication safety cases
& High
& Healthcare professionals need reliable,
organized information to review cases.
& Review prescriptions and medication
information manually.
& \PAChange{Provide structured safety cases,
a simulated HCP review workflow,
and authorized patient--HCP messages.}
\\ \hline

% --------------------------------------------------
% MEDICATION REMINDERS AND ADHERENCE
% --------------------------------------------------

Follow confirmed medication schedules
& High
& Users may forget doses or misunderstand
medication schedules.
& Use paper notes, alarms, or reminder apps.
& Provide reminders based on confirmed
schedules, apply a safety gate to affected
new reminders, and track adherence.
\\ \hline

% --------------------------------------------------
% CONTEXT-AWARE NLP CHATBOT
% --------------------------------------------------

\PAChange{Prepare patient health summaries}
& \PAChange{High}
& \PAChange{Patients may forget important
health and medication information.}
& \PAChange{Collect records manually
or rely on memory.}
& \PAChange{Generate structured summaries
for patient review before sharing.}
\\ \hline

\PAChange{Prepare for medical consultations (Planned)}
& \PAChange{Medium}
& \PAChange{Patients may forget important
concerns or questions during appointments.}
& \PAChange{Prepare personal notes
or rely on memory.}
& \PAChange{Future: help patients organize symptoms,
medications, and questions for doctors.}
\\ \hline

\PAChange{Compare medication records (Planned)}
& \PAChange{Medium}
& \PAChange{Prescriptions may contain
different medications, strengths,
or usage instructions.}
& \PAChange{Compare prescriptions manually.}
& \PAChange{Future: highlight recorded differences
without assuming treatment changes.}
\\ \hline

\PAChange{Understand medication instructions (Planned)}
& \PAChange{Medium}
& \PAChange{Patients may misunderstand
instructions or struggle to follow schedules.}
& \PAChange{Read instructions or
consult healthcare professionals.}
& \PAChange{Future: explain confirmed instructions
and help identify adherence problems
without changing treatment.}
\\ \hline

% --------------------------------------------------
% DATASET REVIEWER
% --------------------------------------------------

\PAChange{Review medication dataset quality}
& \PAChange{High}
& \PAChange{Incorrect medication records
or ingredient mappings may affect
safety assessments.}
& \PAChange{Inspect datasets and
supporting evidence manually.}
& \PAChange{Support dataset review,
record decisions, and track changes
with source references.}
\\ \hline

% --------------------------------------------------
% ADMINISTRATOR
% --------------------------------------------------

\PAChange{Manage role-based access}
& \PAChange{High}
& \PAChange{Different user roles
require different access permissions.}
& \PAChange{Configure accounts and
permissions manually.}
& \PAChange{Allow administrators
to manage accounts, roles,
and permissions.}
\\ \hline

% --------------------------------------------------
% TECHNICAL SUPPORT
% --------------------------------------------------

\PAChange{Resolve technical issues}
& \PAChange{Medium}
& \PAChange{Users may encounter
application errors or access problems.}
& \PAChange{Contact developers
or report issues informally.}
& \PAChange{Support issue reporting,
troubleshooting, and resolution tracking.}
\\ \hline

% --------------------------------------------------
% USABILITY AND SECURITY
% --------------------------------------------------

Use the system easily
& High
& Older adults and users with limited digital
skills may have difficulty using applications.
& \PAChange{Use simple tools or ask
others for help.}
& Provide clear interfaces, readable text,
simple guidance, and alternative input methods.
\\ \hline

Protect personal health information
& High
& Medication records, prescription images,
and voice data may contain sensitive information.
& Store information in personal notes
or separate applications.
& Limit data collection, restrict access,
and use synthetic or non-identifying
test data in the prototype.
\PAChange{Enforce role-based permissions
and prevent unauthorized information sharing.}
\\ \hline

\end{longtable}




% ==================================================
% 3.5 ALTERNATIVES AND COMPETITION
% ==================================================

\subsection{Alternatives and Competition}

\PAChange{The alternatives below include traditional
methods and existing products. Comparisons focus on
publicly described capabilities and the proposed
scope of our prototype, not measured superiority.}

\begin{longtable}{|
    L{0.23\textwidth}|
    L{0.23\textwidth}|
    L{0.24\textwidth}|
    L{0.20\textwidth}|
}
\hline
\textbf{Alternative} &
\textbf{Strengths} &
\textbf{Weaknesses / Limits} &
\textbf{Our Differentiation}\\
\hline
\endfirsthead
\hline
\textbf{Alternative} &
\textbf{Strengths} &
\textbf{Weaknesses / Limits} &
\textbf{Our Differentiation}\\
\hline
\endhead

Direct consultation with a doctor or pharmacist
& Provides qualified medical judgment and
personalized advice.
& May require time, travel, or additional cost.
Patients may provide incomplete medication information.
& Helps users prepare structured medication
summaries and identify cases needing professional review.
It does not replace medical consultation.
\\ \hline

General Web search or AI chatbot
& Easy to access and supports flexible
information searches.
& Information quality and coverage may vary.
AI-generated answers may contain unsupported claims.
& \PAChange{Combines supported medicine references,
explicit uncertainty, and authorized patient records
within a medication-focused workflow.}
\\ \hline

\PAChange{Drugs.com Drug Interaction Checker}
& \PAChange{Checks reported interactions across
multiple medicines and provides severity information.}
& \PAChange{Support for each Vietnamese brand name and
our prescription-confirmation workflow cannot be assumed.}
& \PAChange{Proposes Vietnamese medicine identification,
user confirmation, and linked reminder safety rules.}
\\ \hline

\PAChange{Medscape Drug Interaction Checker}
& \PAChange{Provides multi-drug interaction information
and clinical severity categories.}
& \PAChange{Designed mainly for drug-reference use;
Vietnamese brand coverage must be checked separately.}
& \PAChange{Proposes a patient-focused intake, risk
explanation, and simulated review workflow.}
\\ \hline

\PAChange{Medisafe medication management app}
& \PAChange{Supports dose reminders, adherence tracking,
and interaction information for supported entries.}
& \PAChange{Its interaction view has language and medicine-
entry restrictions; features vary by configuration.}
& \PAChange{Proposes confirmed prescription schedules
and holds new reminders for detected severe interactions.}
\\ \hline

Paper prescriptions, notes, or phone alarms
& Familiar, inexpensive, and can work offline.
& Information may be scattered, difficult to update,
and lack medication safety checks.
& \PAChange{Aims to combine confirmed medication records,
safety findings, and reminders in one Web application.}
\\ \hline


\end{longtable}

% Sources for competitor comparison
\par\smallskip

\noindent\PAChange{
\textit{Sources checked on 01 Nov 2026:}
\href{https://www.drugs.com/drug_interactions.html}
{\underline{Drugs.com Interaction Checker}};
\href{https://reference.medscape.com/drug-interactionchecker}
{\underline{Medscape Interaction Checker}};
\href{https://help-center.medisafe.com/en/articles/8097955-how-can-i-see-possible-med-interactions}
{\underline{Medisafe Help Center}}.
Feature coverage may change and has not been
independently tested by the project team.
}
